\section{Thảo luận}

Trong quá trình thực hiện dự án, nhóm đã gặp một số điểm đáng lưu tâm cần được thảo luận và cải thiện trong các giai đoạn tiếp theo.

\subsection{Những khó khăn gặp phải}

\begin{itemize}
    \item \textbf{Làm sạch dữ liệu:} Tập dữ liệu ban đầu có một số giá trị ngày tháng không đồng nhất, yêu cầu phải chuẩn hóa định dạng trước khi phân tích thời gian.
    \item \textbf{Lựa chọn biểu đồ:} Việc quyết định chọn biểu đồ nào để thể hiện các mối quan hệ phức tạp (như Heatmap hay Scatter plot cho sự kết hợp giữa địa lý và doanh thu) mất khá nhiều thời gian để tối ưu hóa.
\end{itemize}

\subsection{Những hạn chế của thiết kế hiện tại}

\begin{itemize}
    \item \textbf{Tương tác chưa tối ưu:} Dashboard hiện tại chủ yếu là các hình ảnh tĩnh, độ tương tác giữa người dùng và dữ liệu chưa cao (như lọc dữ liệu theo thời gian thực).
    \item \textbf{Thiếu dữ liệu chi tiết:} Một số biến số như "Phản hồi khách hàng" chưa đủ chi tiết để phân tích sâu hơn về lý do trả hàng.
\end{itemize}

\subsection{Hướng phát triển}

\begin{itemize}
    \item \textbf{Triển khai Dashboard tương tác:} Chuyển đổi các thiết kế tĩnh sang các công cụ như Streamlit, Dash hoặc Tableau để người dùng có thể lọc và tùy chỉnh chế độ xem.
    \item \textbf{Tích hợp mô hình dự báo:} Sử dụng Machine Learning để dự đoán doanh thu trong các tháng tới dựa trên xu hướng hiện có.
    \item \textbf{Mở rộng dữ liệu:} Kết hợp thêm các nguồn dữ liệu từ các nền tảng mạng xã hội để phân tích thêm về xu hướng thị trường.
\end{itemize}
